\scanpage{74}
\mspage{II.22}
% archive sheet number: 73
% line 1
\[
H^1(\bar X,\unsure{\mathbb Z_\ell|Z_\ell})^{\pi_0}\subset H^1(\bar X',\mathbb Z_\ell)^{\pi_0}.\tag*{(2.5.4).}
\]
% lines 2–3
\deleted{Il reste à \deleted{\unsure{Démontrer}} \unsure{voir}} que le dernier membre est nul :
\[
H^1(\bar X',\mathbb Z_\ell)^{\pi_0}=1.\tag*{(2.5.5)}
\]
% lines 4–9
C'est une conséquence bien connue de l'hyp. de Riemann-Weil, prouvée en dim. 1. Explicitons l'argument. \deleted{Soit $A$ \deleted{\unsure{la}} \unsure{la jacobienne},} \inserted{Posons $X'=C_k$.} On sait que $\unsure{k|h}$ est le \deleted{\unsure{anneau}} \inserted{corps} des fractions d'une \deleted{\unsure{algèbre}} Algèbre \inserted{$R$} de t.f. sur $\mathbb Z$, \deleted{\unsure{donc}} telle que \deleted{\unsure{X}} $C_k$ provienne d'un schéma $C_S$ sur $S=\operatorname{Spec}R$, propre et lisse sur $S$. \inserted{On peut supposer $S$ \uline{normal}.} La \deleted{homomorphisme} \inserted{relation} (2.5.5) s'écrit alors
% line 10
\[
H^0(S,R^1f_*(\mathbb Z_\ell))=0
\]
% lines 11–16
$f:C-S$. Sous cette forme, on voit que cette relation \deleted{\unsure{entraîne}} ne change pas si on \deleted{change} \inserted{remplace} le pt géom $\unsure{\xi|\eta}$ de $S$ par \deleted{\unsure{parallèle}} \inserted{un} autre $\xi'$, \inserted{en faisant opérer $\pi_1(S,\xi')$ sur $H^1(C_{\xi'},\mathbb Z_\ell)$,} \deleted{\unsure{qu'on}} peut supposer \inserted{\unsure{le choisir en un|le choix en un}} \deleted{fini} \uline{fermé} \inserted{$s$} de $S$, donc à corps résiduel fini. \deleted{Alors} À fortiori, il suffit de prouver la relation
% line 17
\[
H^1(C_{\bar s},\mathbb Z_\ell)^{\pi_1(s,\bar s)}=0\tag*{(2.5.6)}
\]
% lines 18–22
Cela nous ramène au cas d'un corps de base \uline{fini}. Mais dans ce cas, on sait \inserted{$\pi_0\simeq\hat{\mathbb Z}$,} \deleted{\unsure{d'après les}} \inserted{et à l'effet \unsure{le|de} prouver, il suffit de voir que \unsure{Ker} $(1-f)$ dans $H^1(C_{\bar s},\mathbb Z_\ell)\simeq T_\ell(A)^\vee$ est nul, i.e Coker $(1-f)$ dans $T_\ell(A)$ fini, i.e. Ker $(1-f)$ dans $T_\ell(A)$ fini, i.e. car $A(k)$ fini, ce qui est trivial.} \deleted{les valeurs propres du générateur de Frobenius, frob de $\hat{\mathbb Z}$ opérant sur $H^1(C_{\bar s},\mathbb Z_\ell)$ sont de valeur absolue $q^{1/2}$ ($q=\operatorname{Card}k$). \inserted{[2,\unsure{\ldots}]} Donc elles sont $\ne1$, \unsure{d'où} (2.5.6).}

% lines left margin beside 8–9, circled and crossed
\margin{\deleted{\unsure{f fait un} \illegible{2} \unsure{pt.} EGA IV \unsure{\ldots}}}
% lines left margin beside 11–19
% margin note written sideways
\margin{\uline{Remarque} 2.6. Nous avons utilisé exactement 2 propriétés de notre corps $K$ : a) \unsure{l'extension} $\subset\bar K$ obtenue en adjoignant les racines $\ell^{\unsure{\nu|v}}$-ièmes de l'unité, pour $\forall\nu$, est infinie b) Pour tout schéma abélien $A$ sur une extension finie $K'$ de $K$, $T_\ell(A)_\pi$ est fini, où $\pi=\operatorname{Gal}(\overline{K'}/K')$. Pour b) il \uline{ne suffit pas} de savoir que $T_\ell(A)^\pi$ est tjrs nul (exemple : $K=\unsure{\mathbb Q_\ell|Q_\ell}$)!}
